% =============================================================================
%  律 D（带宽自适应律）与 T 自适应律：MathType 可直接读取的 LaTeX 片段
%  版本：v5.0 · 2026-11-12
%
%  用法（MathType）：
%    1) MathType → Preferences → Workspace Preferences
%       勾选 "Allow TeX language input"（允许 TeX 语言输入）；
%    2) 复制下面任意一行的公式（不含行首的 % 注释行）；
%    3) 在 MathType 编辑窗口中粘贴并回车，即转成公式。
%
%  约定：每条公式都是**单行、独立、可单独粘贴**的 LaTeX；
%        不使用 \begin/\end、\label、\tag、\text、\operatorname 与自定义宏，
%        以最大化 MathType 的兼容性。
%        若某台 MathType 不认 \left\| ... \right\|，可改成 \| ... \|。
% =============================================================================


% -----------------------------------------------------------------------------
% A. 记号（对象 / 观测误差 / 内模）
% -----------------------------------------------------------------------------
% A1 被控对象（标幺化）
\dot{y}=b\,u+d,\quad b_0=1.2\,b

% A2 ESO 观测误差
e_1=y-z_1

% A3 低通重复通道（内模）
\tau_0\,\dot{v}+v=Q_0\,v(t-T)+\beta_3\tanh\left(\frac{e_1}{\zeta}\right)


% -----------------------------------------------------------------------------
% B. 律 D：带宽自适应律（k_p 与 omega_o 联合调度）
% -----------------------------------------------------------------------------
% B1 检测量：跟踪误差的 T 窗均值
\bar{e}_1(t)=\frac{1}{T}\int_{t-T}^{t}e_1(\tau)\,\mathrm{d}\tau

% B2 门控（非周期事件判别器）
g(t)=\frac{\bar{e}_1^2(t)}{\bar{e}_1^2(t)+\varepsilon_g^2}

% B3 比例增益调度
k_p^{\mathrm{eff}}(t)=k_p\left[1+k_b\,g(t)\right]

% B4 观测带宽调度
\omega_o^{\mathrm{eff}}(t)=\omega_o\left[1+a\,g(t)\right]

% B5 二阶带宽一致的配套缩放
\beta_1^{\mathrm{eff}}=2\,\omega_o^{\mathrm{eff}},\quad \beta_2^{\mathrm{eff}}=\left(\omega_o^{\mathrm{eff}}\right)^2

% B6 调度后的 ESO
\dot{z}_1=z_2+b_0\,u+v+\beta_1^{\mathrm{eff}}\,G(e_1),\quad \dot{z}_2=\beta_2^{\mathrm{eff}}\,G(e_1)

% B7 类 fal 校正增益（z1 通道指数 1、z2 通道指数 a）
G(e;a,\delta,\lambda)=\left(\left|e\right|+\delta\,\lambda^{-1/a}\right)^{a}\tanh\left(\frac{\lambda e}{\delta}\right)

% B8 控制律
u=\frac{1}{b_0}\left[k_p^{\mathrm{eff}}\,(r-z_1)-z_2-v\right]

% B9 控制律全写（把 B3 代入 B8）
u=\frac{1}{b_0}\left[k_p\left(1+k_b\,g\right)(r-z_1)-z_2-v\right]

% B10 定稿参数（文字，供核对，不必粘贴）
%     k_b=6,  a=3,  epsilon_g=0.01,  beta_3 不调度（b=0）


% -----------------------------------------------------------------------------
% C. T 自适应律（标称周期 + 锁相极值搜索 ESC）
% -----------------------------------------------------------------------------
% C1 内模传递函数
\frac{V(s)}{W(s)}=\frac{1}{1+\tau s-Q\,e^{-sT}}

% C2 分母在 j omega 处的模平方（theta = omega T）
\left|D(j\omega)\right|^2=1+(\omega\tau)^2+Q^2-2Q\left[\cos\theta-\omega\tau\sin\theta\right]

% C3 三角配方（varphi = arctan(omega tau)）
\cos\theta-\omega\tau\sin\theta=\sqrt{1+(\omega\tau)^2}\,\cos\left(\theta+\varphi\right),\quad \varphi=\arctan(\omega\tau)

% C4 相位对齐条件（|D| 的唯一极小点）
\omega T=2\pi k-\arctan(\omega\tau)

% C5 最优周期修正
\delta^{*}(\omega)=-\frac{\arctan(\omega\tau)}{2\pi}

% C6 相位对齐后的最小残余
\left|D\right|_{\min}=\sqrt{1+(\omega\tau)^2}-Q

% C7 两层周期（内层给频率，外层修相位）
T=T_{\mathrm{gross}}\,(1+\delta),\quad T_{\mathrm{gross}}=\frac{2\pi}{\omega_f}

% C8 搜索抖动
\delta(t)=\hat{\delta}+\Delta\sin(\omega_d t)

% C9 基波残余幅值（带通后的一周期均方根）
A[k]=\sqrt{\frac{1}{N_T}\sum_{i=0}^{N_T-1}y_{\mathrm{bp}}^2[k-i]},\quad N_T=\frac{T_{\mathrm{gross}}}{\Delta t_s}

% C10 锁相解调
I_k=\frac{2}{N P_d}\int_{(k-1)NP_d}^{kNP_d}A(t)\sin(\omega_d t)\,\mathrm{d}t

% C11 极值搜索更新律
\hat{\delta}_{k+1}=\mathrm{clip}\left(\hat{\delta}_k-\sigma\,\mathrm{sgn}(I_k),\,-\delta_{\max},\,\delta_{\max}\right)

% C12 目标函数的碗形近似
A(\delta)\approx A^{*}\left[1+\kappa(\delta-\delta^{*})^2\right]

% C13 解调量的线性化
I_k\approx 2\kappa A^{*}\left(\hat{\delta}_k-\delta^{*}\right)\Delta

% C14 可分辨条件（单步调制量大于 A 的噪声）
2\kappa A^{*}\sigma\left|\hat{\delta}-\delta^{*}\right|\ge\eta_A

% C15 分数阶延迟（线性插值）
v[k-n_d]\approx(1-\lambda)\,v[k-i_0]+\lambda\,v[k-i_0-1],\quad n_d=\frac{T}{\Delta t_s},\ i_0=\lfloor n_d\rfloor,\ \lambda=n_d-i_0

% C16 换速冻结判据
\left|\frac{T_{\mathrm{gross}}[k]-T_{\mathrm{gross}}[k-1]}{T_{\mathrm{gross}}[k-1]}\right|>0.05\quad \mathrm{or}\quad \left|1-y\right|>0.1

% C17 换速前馈重置
\hat{\delta}\leftarrow\mathrm{clip}\left(-\frac{\arctan(\omega_f\tau)}{2\pi},\,-\delta_{\max},\,\delta_{\max}\right)

% C18 终版参数（文字）：Delta=0.5%, sigma=1%, P_d=1.5 s, N=3, delta_0=-14.3%, delta_max=20%


% -----------------------------------------------------------------------------
% D. 证明用公式
% -----------------------------------------------------------------------------
% D1 盲性：周期分量在 T 窗均值为零，故检测量只含非周期部分
e_1=p+q,\quad p(t+T)=p(t),\quad \int_0^{T}p(\tau)\,\mathrm{d}\tau=0\ \Rightarrow\ \bar{e}_1(t)=\bar{q}(t)

% D2 门控的取值范围
0\le g<1,\quad g=0\ \Leftrightarrow\ \bar{e}_1=0

% D3 门控的 Lipschitz 常数
\left|g'(x)\right|\le\frac{3\sqrt{3}}{8\,\varepsilon_g}

% D4 调度增益的下界（UUB 证明的起点）
k_p\le k_p^{\mathrm{eff}}<k_p\left(1+k_b\right),\quad \omega_o\le\omega_o^{\mathrm{eff}}<\omega_o\left(1+a\right)

% D5 跟踪误差动态
\dot{\sigma}=-\frac{b}{b_0}\,k_p^{\mathrm{eff}}\,\sigma+\frac{b}{b_0}\,k_p^{\mathrm{eff}}\,e_1+w(t),\quad \sigma=r-y

% D6 跟踪误差的 UUB 界
\limsup_{t\to\infty}\left|\sigma(t)\right|\le\bar{e}_1+\frac{b_0}{b\,k_p}\left(\sup\left|d\right|+\frac{b}{b_0}\left(\sup\left|x_2\right|+\bar{v}\right)\right)

% D6b 观测误差地板（跟踪界的第一项）
\bar{e}_1=\frac{2\sqrt{3}}{\omega_o}\sqrt{\bar{v}^2+\left(\frac{H}{\omega_o}\right)^2}

% D7 重复通道自持界（与 e1 无关）
\limsup_{t\to\infty}\left|v(t)\right|\le\bar{v}=\frac{\beta_3}{1-Q_0}

% D8 观测器误差的 UUB 界
\limsup_{t\to\infty}\left\|\xi(t)\right\|\le\frac{2\sqrt{3}}{\omega_o^{\mathrm{eff}}}\sqrt{\bar{v}^2+\left(\frac{H}{\omega_o^{\mathrm{eff}}}\right)^2}

% D9 时变延迟算子的 L2 增益
\left\|\mathcal{D}_T v\right\|_2^2\le\frac{1}{1-\mu}\left\|v\right\|_2^2,\quad \left|\dot{T}\right|\le\mu<1

% D10 内模的输入-状态稳定性（小增益）
\left\|v\right\|_2\le\frac{1}{1-\gamma}\left\|w\right\|_2,\quad \gamma=\frac{Q}{\sqrt{1-\mu}}

% D11 小增益条件等价形式
\gamma<1\ \Leftrightarrow\ \mu<1-Q^2

% D12 周期变化率的显式表达
\left|\dot{T}\right|=\left|\frac{2\pi}{\omega_f}(1+\delta)\right|\left|\frac{\dot{\omega}_f}{\omega_f}+\frac{\dot{\delta}}{1+\delta}\right|\le\mu<1-Q^2

% D13 测速滤波的设计约束
\tau_f\ge\frac{T\left|\dot{\omega}\right|}{\omega\,\mu_{\max}},\quad \mu_{\max}=1-Q^2

% D14 带宽调度的平均驻留时间条件
\tau_a>\frac{\ln(1+\beta)^2}{\omega_o}

% D15 相邻层 Lyapunov 函数的比值界
\frac{V_{\omega_1}(e)}{V_{\omega_2}(e)}\le\left(\frac{\omega_{\max}}{\omega_{\min}}\right)^2

% D16 名义闭环（调度关）与调度闭环的界一致
k_p^{\mathrm{eff}}\ge k_p,\ \omega_o^{\mathrm{eff}}\ge\omega_o\ \Rightarrow\ \bar{\sigma}_{\mathrm{sched}}\le\bar{\sigma}_{\mathrm{nominal}}

% D17 换速瞬态的解连续依赖性
\left\|v-v_{\mathrm{fr}}\right\|_{\infty}\le C\,\Delta T\,e^{L\delta_t}\left(1+\left\|w\right\|_{\infty}\right)
